\section{Day 11: Bloom Filters (Oct 06, 2025)}

\subsection{Bloom Filters}

A \textbf{bloom filter} is a data structure that approximately maintains a set of keys $S$ from universe $U$. For a given $x$, it will return ``$x$ is definitely not in $S$'' or ``$x$ is probably in $S$''.
\begin{description}
    \item[BF-Insert($S$, $x$)] adds key $x$ to $S$
    \item[BF-Search($S$, $x$)] If $x \in S$, returns $ProbTrue$. If $x \not \in S$, returns False 99\% of the time, but returns $ProbTrue$ with low probability (1\%)
\end{description}

\begin{definition}[False Positive]
When \textsc{Search} returns $ProbTrue$, yet $x \not \in S$.
\end{definition}

The upside is that bloom filters are space efficient. Here, more space means a smaller probability of error. An example usecase is when a large blacklist of websites needs to be stored. A bloom filter consists of
\begin{itemize}
    \item An array $BF[0, \dots, m-1]$ of $m$ bits, initialized to 0 
    \item $h_1, \cdots, h_t$, which are $t$ independent hash functions $U \to \{ 0, \cdots, m-1 \}$. \\
    Assume random hash functions mapping each key in $U$ to $\{ 0, \cdots, m-1 \}$ uniformly at random (for simplicity).
\end{itemize}

\subsection*{Insert($S$, $x$)}
\begin{algorithmic}[1]
    \For{$i = 1$ to $t$}
    \State $BF[H_i(x)] = 1$
\EndFor
\end{algorithmic}

\subsection*{Search($S$, $x$)}
\begin{algorithmic}[1]
    \For{$i = 1$ to $t$}
    \If{$BF[h_i(x)] = 0$}
        \State return $False$
    \EndIf
\EndFor
\State return $ProbTrue$
\end{algorithmic}

\noindent Suppose $x_1, \cdots, x_n$ are inserted into $S$. 
\begin{problem}
    What is the probability that $BF(\ell) = 0$, for a fixed $0 \leq \ell \leq m - 1$? 
\end{problem}

\begin{problem}
    What is the probability of \textsc{BF-Search} returning $ProbTrue$ for $y \not \in S$? (Probability of a getting a false positive)
\end{problem}

\begin{problem}
    Given $n$ and $m$, how to choose $t$ to minimize probability of false positives ($f$). 
\end{problem}

We can tackle this by getting more hash functions, meaning there is a larger chance for \textsc{Search} to find a 0 bit and returning $False$, by having a larger fraction of 1s in the BF array. We do this while making the BF array longer (increasing the size of $m$)?

We can approximate $p$ with $\tilde{p}$, where
\[
p = \left(1 - \frac{1}{m}\right)^{nt} \approx e^{\frac{-tn}{m}} = \tilde{p}
\]
From calculus, the optimal choice of $t$ occurs when $t = \frac{m}{n} \ln (2)$, have $\tilde{{p}} = 0.5$.

\subsection{Open Addressing}

Open addressing is an alternative way to handle collisions. We currently handle collisions by storing them in a doubly-linked list, which is not very efficient.

We store keys directly in the table, and probe a sequence of slots to find each key.
\[
h(k, i) : U \times \{ 0, \cdots, m-1 \} \to \{ 0, \cdots, m-1 \}
\]
The sequence $(h(k, 0), h(k, 1), \cdots, h(k, m-1))$ is called the \textbf{probe sequence}. It must be a permutation of $\{ 0, \cdots, m-1 \}$ (this isn't a strict requirement, but otherwise you won't be checking all of the slots).

More probing methods include linear probing, quadratic probing, and double hashing.

\begin{description}
    \item[Insert($T$, $k$)] read entries of $T$ according to prove sequence, place $K$ in the first empty slot. Return an error if $T$ is full.
    \item[Search($T$, $k$)] traverse the probe sequence looking for $k$. Return $False$ if we find an empty slot, or search all $m$ slots without finding $k$.
    \item[Delete($T$, $k$)] \textbf{lazydeletion or `tombstone method'}: replaced keys with the ``DELETED'' marker.\\
    For \textsc{Insert}, we can insert into a marked slot. However, for \textsc{Search}, we cannot stop if we encounter a marker. The key we're searching for may still be in the probe sequence, thus we cannot stop early.
\end{description}
